\chapter{基于事件的 SLAM}

\textbf{作者：} Guillermo Gallego、Javier Hidalgo-Carrió 和 Davide Scaramuzza

读者不难注意到，在需要解释空间关系并与环境交互的现实应用中，SLAM 至关重要。SLAM 所使用的主要传感器决定了系统的成功与适应能力。视觉 SLAM 是最普及的 SLAM 方法类别之一，因为相机易于获得、价格低廉，并能产生直观且信息丰富的信号，使机器人能够在广泛场景中感知世界（例如构建无需 GNSS 等额外基础设施的轻量系统）。尽管相关技术已经取得进展，当前最先进的人工视觉系统在现实任务中的鲁棒性和效率仍不及生物视觉系统。标准相机以与场景动态无关的固定帧率感知世界，因此在两帧之间的时间段内处于“失明”状态，引入延迟，可能丢失跟踪；当场景中没有运动时，还会产生大量冗余数据。本章探索一个具有前瞻性的挑战：借助受生物启发的硅视网膜或“事件相机”技术，理解并构建快速（不受帧率限制）、低功耗且能适应宽范围照明条件的视觉 SLAM 系统。事件相机克服了标准相机的若干局限（见图 10.1）。

我们首先介绍事件相机的工作原理（第 10.1 节）以及相应的挑战与应用（第 10.2 节），随后聚焦于事件相机数据处理方法（第 10.3 节）以及相应的前端（第 10.4 节）和后端处理（第 10.5 节）。最后讨论最先进的系统（第 10.6 节）、数据集、模拟器与基准（第 10.7 节），并在第 10.8 节以新趋势和延伸阅读作结。

\section{传感器描述}

\subsection{工作原理}

与由外部时钟（例如 30 Hz）决定采集整幅图像速率的传统相机不同，事件相机（例如 Dynamic Vision Sensor，DVS）

\begin{figure}[htbp]
\centering
\includegraphics[width=\linewidth,height=0.78\textheight,keepaspectratio]{91588c3e5554144d7bbe083c816b89e4fde0c9d993f20d04a7c307245a6e036e.jpg}
\caption{图 10.1 配备下视 DAVIS 相机 \textsuperscript{[111]}（240×180 像素）的无人机，使用视觉惯性里程计（VIO）算法 \textsuperscript{[944]} 进行状态估计并完成自主飞行。事件相机数据的高速和高动态范围特性使系统能够在困难照明条件下运行。插图展示了在灰度帧（左，存在运动模糊）和事件扭曲后经运动补偿的图像（中，清晰）中检测并跟踪的特征（即关键点）。事件数据根据极性以红/蓝色显示，清楚地响应场景轮廓。同一 VIO 算法 \textsuperscript{[944]} 还在高速场景中得到演示，例如将事件相机系在绳子上旋转。图像来自 \textsuperscript{[361]}（©2020 IEEE）。}
\label{fig:chapter-10-91588c3e5554144d7bbe083c816b89e4fde0c9d993f20d04a7c307245a6e036e}
\end{figure}

的像素彼此独立工作，异步地响应场景中发生的亮度变化（图 10.2b）。逐像素变化由场景照明（例如闪烁的灯光）和/或相机与场景（包括运动物体）之间的相对运动引起。因此，事件相机的输出是一系列数字“事件”（或“脉冲”），每个事件表示亮度（对数强度）的变化。这种编码受生物视觉通路的脉冲特性启发（图 10.2a）。

具体而言，每个像素在发送事件时记忆对数强度 $L$，并持续监测相对于该记忆值、幅度足够大的变化 $\Delta L$（图 10.2）。当变化达到阈值 $C$ 时，

\begin{equation}
\Delta L \doteq L (\boldsymbol {x} _ {k}, t _ {k}) - L (\boldsymbol {x} _ {k}, t _ {k} - \Delta t _ {k}) = p _ {k} C,\tag{10.1}
\end{equation}

相机便发送事件 $e _ { k } \doteq ( \boldsymbol { x } _ { k } , t _ { k } , p _ { k } )$。该事件从芯片传出，其中包含 $x,y$ 像素位置 $\scriptstyle { \mathbf { { \mathit { x } } } } _ { k }$、时间 $t _ { k }$ 以及变化的 1 位极性 $p _ { k } \in \{ + 1 , - 1 \}$（即亮度增加或减小）。$\Delta t _ { k }$ 是同一像素上前一事件以来经过的时间。

事件相机是数据驱动传感器：其输出取决于场景中的运动量或照明变化量。运动越快，每秒产生的事件越多，因为每个像素会根据所监测强度信号的变化速率自适应其采样率。

\begin{figure}[htbp]
\centering
\includegraphics[width=\linewidth,height=0.78\textheight,keepaspectratio]{942aafbf56ce3bb9e2cb412fca8648b0c57540084d1965a836e542ddcf415c74.jpg}
\caption{图 10.2 事件相机（例如 DVS）的工作原理：(a) 人类视网膜三层模型及相应的 DVS 像素电路；(b) DVS 像素的工作示意图，将光转换为事件（脉冲），信号颜色与 (a) 中各层相匹配；(c) 标准相机和事件相机对旋转圆盘上的黑点这一视觉刺激的响应对比。事件相机连续传输亮度变化，在时空中形成事件螺旋。红色表示正事件（ON 脉冲），蓝色表示负事件（OFF 脉冲）。图像改编自 \textsuperscript{[884]}（©2014 IEEE）。}
\label{fig:chapter-10-942aafbf56ce3bb9e2cb412fca8648b0c57540084d1965a836e542ddcf415c74}
\end{figure}

\textbf{生物启发：瞬态通路。} 事件相机的设计受到生物视觉通路工作方式的启发；视觉通路是动物和人类处理信息的路径。根据双流假说，背侧通路（也称“瞬态”或“where”通路）专门处理动态视觉信息（例如场景中的运动），而腹侧通路（称为“持续”或“what”通路）专门负责物体和视觉目标的识别。DVS \textsuperscript{[661]} 对应从光感受器到神经节细胞的瞬态通路部分，采用简化的三层像素设计，在生物逼真度与电路稳定性之间取得平衡（图 10.2）。三层分别实现光转换、增量调制和比较功能。Asynchronous time-based image sensor（ATIS）\textsuperscript{[883]} 或 Dynamic and Active-Pixel Vision Sensor（DAVIS）\textsuperscript{[111]} 等相机同时建模两条视觉通路，因此输出两类信号：DVS 事件和灰度信息（例如图像）。关于主要事件相机类型的更多细节见 \textsuperscript{[884, 361]}。

\subsection{事件相机的优势}

事件相机的感知原理与过去七十年主导计算机视觉和机器人视觉的标准（基于曝光）相机截然不同，并带来诸多优势：

\textbf{高时间分辨率：} 事件以微秒级分辨率检测并加盖时间戳，因此能够捕获极快运动，而不会受到帧式相机常见的运动模糊影响。事件几乎连续地产生，从而避免标准相机中可能造成帧间大位移并破坏数据关联的盲时间。

\textbf{低延迟：} 每个像素独立工作，无需等待全局曝光时间；一旦检测到亮度变化，事件便立即传输，延迟低于 1 ms。

\textbf{低功耗和低带宽：} 事件表示非冗余的时间数据，功耗仅用于有用信息。与以相同速率运行的传统相机相比，带宽也得到降低。在芯片层面，相机功耗低于 10 mW，使嵌入式系统的功耗可低至 100 mW 或更低 \textsuperscript{[31]}。

\textbf{高动态范围（HDR）：} 事件相机可感知的光值范围非常大（通常超过 120 dB，而标准相机约为 60 dB），能够同时感知极暗（月光）和极亮（日光）区域。因此不会出现帧式相机常见的欠曝光或过曝光。这一特性源于两个事实：每个像素独立工作，并以对数尺度将光转换为电压。

\subsection{当前器件与趋势}

“我应该购买或使用哪种事件相机来解决 SLAM 问题？”刚进入这一新兴领域的人经常提出这个问题。事件相机的特性通常通过表格进行比较 [361, 表 1]、[162, 表 1–2]。虽然存在多种事件相机设计，但大多数仍是实验室原型；只有少数实现了商业化，从而支持针对经典问题和新问题（如基于事件的 SLAM）探索新方案。在主要制造商（SONY、Samsung、iniVation / SynSense、Prophesee、Omnivision）商业化的器件中，有几项趋势值得关注：

\textbf{像素尺寸：} 随着传统相机的百万像素竞赛和工业需求压力，像素间距（即尺寸）已从 40 µm（DVS128 \textsuperscript{[661]}）显著减小到小于 5 µm \textsuperscript{[326]}。DVS 像素比传统像素执行更多操作（调制、比较等），因而需要更多晶体管，难以在相同传感器面积内封装。为最大化像素受光面积（即填充因子）并缩小像素感光部分之间的间隙，已经采用堆叠技术和背照式技术 \textsuperscript{[326]}。

\textbf{灰度输出：} DAVIS 或 ATIS 等早期器件同时输出灰度数据（例如图像 \textsuperscript{[111]}），这对固定相机应用尤其有用（尽管这并非 SLAM 的通常场景）。受百万像素竞赛驱动，HD 事件相机 \textsuperscript{[326]} 等新型号停止灰度输出，以便为事件输出留出更多面积。

颜色对许多运动相关任务并非必需，因此只有少数事件相机型号提供彩色滤光片，用于检测各颜色通道（红、绿、蓝，即 RGB）的变化 \textsuperscript{[770]}。

\textbf{惯性数据：} 一些相机还提供集成在同一器件中的惯性测量单元（IMU）数据。IMU 是相机有价值的互补本体传感器，可实现视觉惯性里程计（VIO）和 SLAM，并比单传感器系统具有更高的鲁棒性和精度。

不能想当然地认为高空间分辨率事件相机必然优于低分辨率相机。虽然捕获精细空间细节很重要，但噪声和带宽同样影响目标应用的需求。在 SLAM 及相关任务中，事件相机可能快速运动和/或经过纹理丰富的场景；HD（1 百万像素）事件相机每秒可产生数亿个事件。这会导致相机输出传输总线饱和以及高处理需求；目前不存在能实时处理如此事件速率的算法与硬件组合（除非转为阵列表示和/或进行子采样/降采样）。人们正在开发新的混合传感器，例如 \textsuperscript{[1238]}（事件空间分辨率低于强度输出），以及模仿生物视觉以降低带宽的中央凹传感器 \textsuperscript{[317]}，它们可能为上述问题提供替代方案。在 SLAM 中，低像素分辨率（例如 QVGA）更适合算法原型开发和计算受限机器人的实时运行。通常，视场（光学系统）的选择与像素数量同样重要。

\section{挑战与应用}

事件相机代表了视觉数据采集领域的革命性技术。因此，需要设计新的方法（算法和硬件）处理所采集的数据并提取有价值的信息，以发挥传感器优势。主要挑战包括：

\textbf{处理时空输出：} 事件相机的输出与标准相机根本不同：事件是异步且空间稀疏的，而图像是同步且稠密的。因此，为图像序列设计的视觉 SLAM 算法不能直接应用于事件数据。

\textbf{处理运动相关数据：} 与图像不同，每个事件仅包含二值（增加/减小）亮度变化信息，该信息不仅取决于场景纹理，还取决于场景与相机之间的相对运动。

\textbf{处理噪声和动态效应：} 事件相机存在噪声，来源包括固有的光子散粒噪声、晶体管电路噪声、对入射光量的依赖、非理想因素以及低功耗（亚阈值）运行。

这些挑战要求重新思考事件数据的时空、光度和随机特性。在 SLAM 中，这引出了如下问题：从事件中提取位姿或深度信息的最佳方式是什么？应采用何种地图和相机轨迹表示，以考虑事件数据近似连续的时间粒度和稀疏性？如何在运动相关数据下建立对应关系（数据关联）？如何建模问题（及其解），而不引入帧式技术的典型瓶颈？

上述问题推动了基于事件的 SLAM 研究（图 10.3）。这一主题既独立发展，也与其他任务结合；例如，基于事件的 SLAM 促进了其他事件视觉任务的研究。早期工作 \textsuperscript{[229, 570]}（旋转运动 SLAM）和 \textsuperscript{[571]}（6 自由度 SLAM）已经利用 SLAM 与图像重建（从事件恢复绝对强度）之间的协同关系。基于事件的 SLAM 与光流估计也在 \textsuperscript{[229, 1243, 1010, 585]} 中联合处理。

\section{基于事件的 SLAM 方法概述与分类}

根据同时处理的事件数量，基于事件的 SLAM 方法大致分为两类：(i) 逐事件运行的方法，系统状态（例如场景地图和相机轨迹）可在单个事件到达时改变，从而实现最低延迟；(ii) 对事件组/批次/切片/数据包运行的方法，会引入一定延迟。后一类的关键设计选择是数据包大小；已有许多方案，例如固定事件数量、固定时间长度和混合判据。

正交地，根据事件的处理方式，还可区分基于模型的方法和数据驱动（即机器学习）方法。仿照帧式 SLAM 的分类，基于事件的 SLAM 可分为间接法（基于特征，使用事件角点、直线、法向流等）和直接法（使用全部事件）。这一分类与所用目标函数或损失函数的类型有关：几何型与光度型（例如事件极性或事件速率/活动度的函数）；也与总体思想有关：间接法通常包含两个步骤（特征提取，将事件“转换”为几何基元；随后运行几何 SLAM 流水线），直接法通常只有一步，将事件数据映射为运动和场景参数。对于直接法，事件生成模型 (10.1)（或其线性化版本 \textsuperscript{[361]}）是设计估计方法的基石。事件间数据关联是事件视觉、尤其是 SLAM 的核心问题。由于事件相机具有高时间分辨率，数据关联通常通过时间和空间邻域处理；硬关联和软关联策略均已得到研究。

\begin{figure}[htbp]
\centering
\includegraphics[width=\linewidth,height=0.78\textheight,keepaspectratio]{b7983bc219141df3389377b2476099342af9addffe2ea6cfcbe86f439d06421d.jpg}
\label{fig:chapter-10-b7983bc219141df3389377b2476099342af9addffe2ea6cfcbe86f439d06421d}
\end{figure}

EVO：实时事件 6 自由度并行跟踪与建图的几何方法

\begin{figure}[htbp]
\centering
\includegraphics[width=\linewidth,height=0.78\textheight,keepaspectratio]{e3fa4b1556a576502a0314644f3c69b71f2c3f7189bb7be70dbf6ba272e15498.jpg}
\label{fig:chapter-10-e3fa4b1556a576502a0314644f3c69b71f2c3f7189bb7be70dbf6ba272e15498}
\end{figure}

EDS：事件辅助直接稀疏里程计

\begin{figure}[htbp]
\centering
\includegraphics[width=\linewidth,height=0.78\textheight,keepaspectratio]{477a06fdcd350a27d6c405e950e11b6d878d21971d30a8a34029e78bea0ae6fe.jpg}
\label{fig:chapter-10-477a06fdcd350a27d6c405e950e11b6d878d21971d30a8a34029e78bea0ae6fe}
\end{figure}

CMax-SLAM：使用对比度最大化的基于事件的旋转运动束调整与 SLAM

\begin{figure}[htbp]
\centering
\includegraphics[width=\linewidth,height=0.78\textheight,keepaspectratio]{3ad912df3bfa2c1bac48f0d4d2cff3aca5450cd138095caeb51c48dd085b30d7.jpg}
\label{fig:chapter-10-3ad912df3bfa2c1bac48f0d4d2cff3aca5450cd138095caeb51c48dd085b30d7}
\end{figure}

事件相机实时三维重建与 6 自由度跟踪（ETAM）

\begin{figure}[htbp]
\centering
\includegraphics[width=\linewidth,height=0.78\textheight,keepaspectratio]{621be325dc394fd53b2995c35daa7fd6e13eecafd6801d3f2e95226f12b19e3b.jpg}
\label{fig:chapter-10-621be325dc394fd53b2995c35daa7fd6e13eecafd6801d3f2e95226f12b19e3b}
\end{figure}

\begin{figure}[htbp]
\centering
\includegraphics[width=\linewidth,height=0.78\textheight,keepaspectratio]{0301f74690d97f7919593d2469513012c6fd38f34e615c0dbf35488021e43c95.jpg}
\label{fig:chapter-10-0301f74690d97f7919593d2469513012c6fd38f34e615c0dbf35488021e43c95}
\end{figure}

DEVO：深度事件视觉里程计

ESVO2：使用双目事件相机的直接视觉惯性里程计
\begin{figure}[htbp]
\centering
\includegraphics[width=\linewidth,height=0.78\textheight,keepaspectratio]{148811a6766bc0b1daedaa9842e64ac64d8fa32438351f6da9dad8f28cc22442.jpg}
\label{fig:chapter-10-148811a6766bc0b1daedaa9842e64ac64d8fa32438351f6da9dad8f28cc22442}
\end{figure}

基于高斯过程连续时间方法的事件双目视觉里程计（GPCT）
图 10.3 基于事件的 SLAM 正在积极研究中，系统涵盖多种方法，包括经典方法和较新的深度学习方案。由于事件由图像平面中运动的边缘触发，因此以边缘形式（例如稀疏或半稠密三维地图）恢复场景地图是自然选择。图像改编自 EVO \textsuperscript{[914]}（©2017 IEEE）、EDS \textsuperscript{[461]}（©2022 IEEE）、CMax-SLAM \textsuperscript{[416]}（©2024 IEEE）、Kim 等 \textsuperscript{[571]}（©2016 Springer）、ESVO2 \textsuperscript{[812]}（©2025 IEEE）、DEVO \textsuperscript{[585]}（©2024 IEEE）和 Wang 等 \textsuperscript{[1154]}（©2023 IEEE）。

上述各类别都有优点和缺点。使用事件相机求解 SLAM 很有挑战性，历史上通常沿多个维度逐步增加复杂度：未知量数量（自由度，即 DoF）、运动类型（从旋转或二维场景到 6 自由度运动）、场景复杂度（纹理）及其运动状态（静态与动态，即独立运动物体，IMO）。基于事件的 SLAM 并非孤立问题：如第 10.2 节所述，它与光流、跟踪、分割等问题存在联系，联系强弱取决于所作假设或具体场景。除上述趋势外，还可观察到早期研究主要采用基于模型的方法，而近期论文开始探索深度学习方法的可能性。

\section{基于事件的 SLAM 系统前端}

基于事件的 SLAM 系统通常由多个模块组成，分别处理特征提取、数据关联、引导初始化、位姿估计、深度估计等较小子问题。系统的一项基本划分是前端与后端。从输入输出角度看，前端接收原始传感器数据（以及相机标定等可能的辅助信息），输出一组事件相机位姿和场景地图（见图 10.4）。后端细化这些变量（即 SLAM 问题的未知量），改善它们与传感器数据之间的拟合。后端在前端之后运行，速率较低（取决于变量数量），并可将输出反馈给前端，以减少漂移和校正错误。

因此，前端将传感器信息（例如光子）转换为几何基元（例如相机位姿），同时也提取光度信息（例如地图外观）。这通常包含“特征”或“信息”提取步骤。首要挑战是理解事件流包含的信息，并使用保留数据特性（低延迟、稀疏性、HDR 等）的方法提取这些信息。在照明恒定的假设下，事件由运动轮廓（边缘）产生。因此，可以把运动中的事件相机视为异步边缘检测器，并以场景轮廓的形式表述 SLAM 问题（图 10.3）。这在先验上是合理的，因为轮廓是图像平面中信息量最大的区域，可据此估计视网膜运动，进而推断三维信息。每个事件只包含一个四元组并带有噪声，携带的信息很少；因此需要大量事件（例如数千或数百万个）才能可靠估计相机位姿和场景地图。事件流中的信息提取取决于任务和诸多设计选择，例如 SLAM 变量（场景地图、相机轨迹）的表示形式、用于处理数据的硬件以及输出速率等。

基于事件的 SLAM
\begin{figure}[htbp]
\centering
\includegraphics[width=\linewidth,height=0.78\textheight,keepaspectratio]{929950bfdc51793cbf0de7c1a7d411a2ae132a68b1ef2d81b62147eed7a65542.jpg}
\caption{图 10.4 基于事件的 SLAM 流水线：前端计算地图和相机位姿，后端细化地图与位姿。由于事件响应运动边缘，恢复的地图通常是边缘（梯度）地图。示例展示了直接旋转 SLAM 流水线（位姿仅由旋转组成，地图简化为全景地图）\textsuperscript{[417]}。还可通过泊松积分恢复绝对强度地图。}
\label{fig:chapter-10-929950bfdc51793cbf0de7c1a7d411a2ae132a68b1ef2d81b62147eed7a65542}
\end{figure}

\subsection{预处理与事件表示}

在 SLAM 问题中，事件相机在穿过场景时连续输出数据。从相机视角看，所有场景边缘似乎都在运动，因此事件会在图像平面“各处”触发。由于事件稀疏且具有微秒级时间分辨率，每个事件都对应不同的相机位姿。这与传统（帧式）相机截然不同：一幅图像中的所有像素测量具有相同时间戳，因而共享同一相机位姿（传统多视图几何 \textsuperscript{[437]} 正是建立在这一范式上）。许多 SLAM 方法将事件数据转换为其他表示（事件图像、时间图或“时间表面”、体素网格等）\textsuperscript{[361]}，原因包括兼容传统计算机视觉方法和便于解释等。这一步通常意味着信息量化（例如将时间戳相近的事件分组）和/或稀疏性损失（例如在无事件的位置用零填充数组）。

因此，事件表示研究 \textsuperscript{[361, 370]} 受到越来越多关注。它通常是前端的第一阶段，并强烈影响后续处理：事件被转换为更熟悉、易处理的表示（例如图像），从而可以输入为传统图像设计的成熟 SLAM 方法，或设计基于图像的学习方法。之所以进行这种转换，部分原因是研究界仍在探索从事件流提取信息的最佳方式，并试图复用成熟的图像方法。前端可以使用不同事件表示；例如 EVO \textsuperscript{[914]} 在建图模块（EMVS \textsuperscript{[916]}）中使用原始事件，在相机跟踪模块中使用事件（类似边缘）图像。理想情况下，应设计使用事件表示的 SLAM 方法，保留事件相机的高速和稀疏特性，避免传统相机的量化时间、延迟和非稀疏性问题。实践中，这是一个需要异步地重新思考视觉处理的新兴研究主题，仍有大量基础问题有待改进和研究。

\subsection{间接法}

前端的设计选择会很大程度影响系统的其余部分。一项主要选择是采用间接法还是直接法。间接法大致包含两步：首先从事件中提取并跟踪基于点、基于直线或其他类型的特征；然后利用经典 SLAM 的结果，根据这些几何基元估计相机运动和三维场景结构。特征将事件数据压缩为少量信息丰富的基元，使计算资源得到更有效和高效的利用。核心问题是建立并维护事件特征（以及地图地标）之间的对应关系，即数据关联。由于每个事件携带的信息很少且与运动有关，无法仅凭单个事件无歧义地确定关联。得益于事件相机的高时间分辨率，可在像素空间中依据时空邻近性建立关联。因此，跟踪特征比匹配特征更自然。

给定场景地图，相机位姿估计或相机跟踪通常表述为特征配准/对齐问题，通过最小化几何目标（例如以像素空间欧氏距离度量的重投影误差）求解。场景三维结构通常通过三角化（即反投影）对应特征位置得到，例如使用给定的相机位姿获得三维点和直线。可以利用成熟的大量几何方法（多视图几何 \textsuperscript{[437]}）。

与常规视觉 SLAM 类似，间接事件法高度依赖鲁棒的特征提取和跟踪。然而，由于要处理大噪声、稀疏性、异步性和运动依赖等独特挑战，这些组件尚不如帧式方法成熟，限制了系统精度和适用性。为应对这些问题，一些系统采用传感器融合（灰度图像和/或 IMU 数据）。

\subsection{直接法}

直接法使用所有可用数据（而非仅使用符合特征定义的事件数据）估计相机运动和三维场景结构。它不显式提取特征，而是直接将事件数据与地图、图像或其他事件对齐。当事件速率相对于系统处理能力过高时，会采用去噪或子采样等数据削减机制，以减少需要处理的事件数量 \textsuperscript{[416, 572]}。

由于直接法只有一个步骤，运动（相机跟踪）或场景参数（建图）通过优化某个目标函数（例如光度误差或空间事件速率误差）获得。基于光度的目标会诱导几何配准目标；问题未知量由事件传达的类边缘亮度模式与对应图像或地图像素的对齐确定。直接法依赖事件数据的近似连续性，例如可以从上一次估计的位姿计算增量相机位姿：事件连续触发，没有间隙或盲时间，因此增量很小。

直接法中，一个因其最先进的精度而受到重视的子类，是通过事件对齐估计运动或场景参数的方法，其结果表现为清晰的事件扭曲图像（IWE）。其思想是通过“抵消运动”来估计运动，即寻找能够对事件数据进行运动补偿的参数。事件对齐可用方差、梯度幅值、离散度等不同目标度量；这类方法也称聚焦最大化或对比度最大化（CMax）\textsuperscript{[360, 1010]}。在事件可被扭曲到少数像素或一条直线的问题中，这些方法可能出现不希望的全局最优解 \textsuperscript{[1009]}。直接法中的数据关联通常由像素网格距离隐式、软性地处理；不过，在某些情况下，使用最近邻值的硬关联同样可行且有效。

\subsection{基于模型与基于学习的方法}

迄今为止，大多数基于事件的 SLAM 方法都是人工设计的，依据事件相机和 SLAM 问题的工作原理及人类直觉构建。相比之下，深度学习方法利用人工神经网络（ANN）对事件数据建模：可以先将事件转换为类似图像的表示，也可以直接用脉冲神经网络（SNN）处理事件。根据监督信号类型，这些方法通常分为有监督和自监督方法。自监督方法利用事件或其他传感器（例如共址灰度相机）中的时间动态或光度一致性损失，估计深度和相机位姿 \textsuperscript{[1243]}；有监督方法需要地面真值数据训练 \textsuperscript{[585]}，而在现实世界中获取此类数据通常很困难。近年来，人们已在汽车、无人机等平台上采集许多多模态数据集，为 ANN 训练提供了所需数据。

基于学习的方案可以替代 SLAM 流水线的一部分，例如特征提取与跟踪 \textsuperscript{[760]}，也可以尝试替代整个系统（端到端）。这类方法能够隐式处理复杂数据表示和噪声，但需要大量训练数据；当应用于与训练集显著不同的场景（即发生“域偏移”）时，还可能出现泛化问题。

\section{基于事件的 SLAM 系统后端}

SLAM 后端等细化模块 \textsuperscript{[136]} 的目标是改善 SLAM 问题变量与传感器数据之间的一致性，从而提高拟合的精度和鲁棒性，减少系统跟踪与建图模块之间的误差传播。如前面章节所述，束调整（BA）通常用于视觉 SLAM 后端。

基于事件的 BA 仍处于起步阶段，因为多数基于事件的 SLAM 系统缺少细化步骤。相反，它们以并行跟踪和建图方式运行 \textsuperscript{[571, 914, 1302]}，每个模块同时依赖另一个模块的输出作为输入才能正常工作。当前事件 SLAM 系统优先考虑简洁性以及事件相机的低延迟优势，而非精度和鲁棒性。除第 10.2 节所述挑战（噪声、运动相关外观等）外，事件后端还需要联合估计相互关联的变量，这会产生高维搜索空间，使优化在复杂度和延迟方面代价高昂，并容易陷入局部极小值。

直到最近，BA 才开始用于包含事件相机的系统。由于 SLAM 后端的设计高度取决于前端输出（要获得最佳性能，两个模块必须紧密集成），因此我们将事件后端分为间接后端（基于特征）和直接后端（基于光度）。

间接后端继承自经典的间接帧式方法 \textsuperscript{[1107, 645, 787]}。它们在事件流中检测到的几何基元（角点、直线等）上运行，也可能先将事件转换为图像 \textsuperscript{[204, 944]}，以复用帧式检测器。其目标通常是最小化重投影误差，即图像平面中的欧氏距离 \textsuperscript{[437]}。这种方法的优点是能够复用经典 SLAM 中成熟且鲁棒的技术；但它丢弃了事件中包含的大量信息（图像重建方法 \textsuperscript{[918, 1280]} 已揭示这些信息），且往往难以达到期望性能：受噪声和事件运动依赖性的影响，当前事件角点不如帧式角点精确稳定，因此在 SLAM 中的应用较少 \textsuperscript{[611]}。间接后端的例子包括 \textsuperscript{[944, 204, 1154]}。

直接后端处理传感器数据（而非几何基元），目标通常是最小化某种形式的光度误差，因此比间接后端更适配事件数据。\textsuperscript{[461]} 等利用共址帧的灰度信息，借用帧式系统的后端 \textsuperscript{[308, 26]}。然而，灰度帧可能出现运动模糊且动态范围较低。仅事件后端不受这些限制；由于发展较晚，目前主要针对受限运动（平面或旋转）。其目标可以是最大化事件对齐（也称运动补偿或 CMax）\textsuperscript{[359, 416]}，或最小化每个事件传达的光度误差（即时间对比度）\textsuperscript{[417, 418]}。这些方法依据事件生成模型 (10.1) 设计。由于单个事件携带的信息很少，而 SLAM 中问题未知量通常很多，精确 BA 需要大量事件，对计算资源、功耗和延迟提出要求。在自然场景和 6 自由度运动下，研究高效的直接、纯事件 BA 仍有很大空间。

\section{最先进的系统}

表 10.1 汇总了基于事件的 VO/SLAM 具体系统，并按照前文分类描述其特性（直接、间接等）。本章不可能也无意详细介绍所有系统，但有几项趋势值得关注。

文献仍由基于模型的系统主导，数据驱动方法尚未取而代之（尽管未来可能发生这种转变，正如其他计算机视觉任务一样）。从研究伊始，事件相机 SLAM 就在不同假设下处理问题，并逐步提高以下方面的复杂度：(i) 相机运动；(ii) 场景类型；(iii) 额外传感器（或信息，例如场景地图），以简化问题。例如，连接到事件相机的深度传感器可减轻仅从事件估计深度的负担，IMU 则可提供准确的角速度信息。

当某种事件方法表现良好后，通常沿着近乎标准的“利用”路线逐步改进（类似帧式 SLAM）：例如，将单目方法 \textsuperscript{[914]} 扩展到双目或多相机场景 \textsuperscript{[381]}；将 \textsuperscript{[1302]}（以及 \textsuperscript{[585]}）等纯事件方法通过融合惯性数据 \textsuperscript{[685, 812]}（以及 \textsuperscript{[408]}）增强鲁棒性；还可扩展基础系统以处理全向镜头等。尽管遵循这种利用路线，基于事件的 SLAM 仍是新兴领域，处于探索不同技术的阶段。分析表 10.1 中的方法即可看出：研究者利用多样的思想和原理，采用不同地图表示、事件表示和损失函数，设计支撑这些系统的估计算法。尤其需要继续研究能够发挥传感器本质特性的状态估计方法。

\section{数据集、模拟器与基准}

基于事件的视觉系统在原型开发、训练和基准测试中，对高质量、多样且信息丰富的真实与合成数据有很高需求。开发模拟器、数据集和排行榜对于推动领域发展、建立共同且坚实的科学技术基础至关重要。下面介绍事件相机领域具有代表性的 SLAM 数据集、基准和模拟器。

（表 10.1 的方法清单沿用原始出版物中的系统名称、缩写、输入类型和运动类别；各列含义见下方表注。）

\scriptsize
\setlength{\tabcolsep}{2pt}
\begin{longtable}{p{0.100\linewidth}p{0.100\linewidth}p{0.100\linewidth}p{0.100\linewidth}p{0.100\linewidth}p{0.100\linewidth}p{0.100\linewidth}p{0.100\linewidth}p{0.100\linewidth}}
\toprule
系统 & 方法类型 & 直接/间接 & 事件表示 & 全局细化（BA） & 运动 & 场景 & 输入 & 备注 \\
\midrule
\endfirsthead
\toprule
系统 & 方法类型 & 直接/间接 & 事件表示 & 全局细化（BA） & 运动 & 场景 & 输入 & 备注 \\
\midrule
\endhead
Cook \textsuperscript{[229]} & 基于模型 & 直接 & Event Frame & 否 & 旋转 & 自然场景 & E & 使用光流的交互网络 \\
Weikersdorfer \textsuperscript{[1175]} & 基于模型 & 间接 & Individual Event & 否 & 平面 & 二维黑白 & E & 首个基于滤波器的 Ev-SLAM \\
PF-SMT \textsuperscript{[570]} & 基于模型 & 直接 & Individual Event & 否 & 旋转 & 自然场景 & E & 两个交错的贝叶斯滤波器 \\
Censi \textsuperscript{[160]} & 基于模型 & 直接 & Event Packet & 否 & 6 自由度 & 黑白 & E+F+D & 基于图像梯度的滤波式 VO \\
EB-SLAM-3D \textsuperscript{[1176]} & 基于模型 & 直接 & Individual Event & 否 & 6 自由度 & 自然场景 & E+D & 使用深度传感器增强事件 \\
Yuan \textsuperscript{[1260]} & 基于模型 & 间接 & Event Frame & 否 & 6 自由度 & 黑白 & E+I+M & 基于竖直线的相机跟踪 \\
Kueng \textsuperscript{[611]} & 基于模型 & 间接 & Local Point Set & 否 & 6 自由度 & 自然场景 & E+F & 基于事件特征跟踪的 VO \\
ETAM \textsuperscript{[571]} & 基于模型 & 直接 & Individual Events & 否 & 6 自由度 & 自然场景 & E & 三个交错的滤波器 \\
CMax-$\omega$ \textsuperscript{[356]} & 基于模型 & 直接 & Individual Events & 否 & 旋转 & 自然场景 & E & 对比度最大化 \\
EVO \textsuperscript{[914]} & 基于模型 & 直接 & Edge Map & 否 & 6 自由度 & 自然场景 & E & 事件--事件几何对齐 \\
EVIO \textsuperscript{[1303]} & 基于模型 & 间接 & Point sets & 否 & 6 自由度 & 自然场景 & E+I & 基于滤波器和运动补偿特征 \\
Rebecq \textsuperscript{[915]} & 基于模型 & 间接 & MC Event Images & 是 & 6 自由度 & 自然场景 & E+I & 基于特征，滑动窗口后端 \\
RTPT \textsuperscript{[925]} & 基于模型 & 直接 & Individual Events & 否 & 旋转 & 自然场景 & E & 全景跟踪与建图 \\
Gallego \textsuperscript{[358]} & 基于模型 & 直接 & Individual Events & 否 & 6 自由度 & 自然场景 & E+M & 鲁棒传感器模型 \\
Mueggler \textsuperscript{[779]} & 基于模型 & 直接 & Individual Events & 是 & 6 自由度 & 自然场景 & E+I+M & 连续时间位姿估计器 \\
USLAM \textsuperscript{[944]} & 基于模型 & 间接 & MC Event Images & 是 & 6 自由度 & 自然场景 & E+F+I & 传感器融合与滑动窗口后端 \\
Chin \textsuperscript{[204]} & 基于模型 & 间接 & Event Frames & 是 & 旋转 & 星场景 & E & 针对星跟踪设计 \\
ESVO \textsuperscript{[1302]} & 基于模型 & 直接 & Time surfaces (TS) & 否 & 6 自由度 & 自然场景 & 2E & 基于 TS 图像块的双目匹配 \\
Hadviger \textsuperscript{[425]} & 基于模型 & 间接 & Corners on TS & 否 & 6 自由度 & 自然场景 & 2E & 交叉相关特征描述子 \\
CMax-GAE \textsuperscript{[572]} & 基于模型 & 直接 & Individual Events & 否 & 旋转 & 自然场景 & E & 对比度最大化 \\
EKLT-VIO \textsuperscript{[725]} & 基于模型 & 间接 & Individual events & 是 & 6 自由度 & 自然场景 & E+F+I & EKLT 跟踪器和 VIO 后端 \\
EDS \textsuperscript{[461]} & 基于模型 & 直接 & Event images & 是 & 6 自由度 & 自然场景 & E+F & 帧式后端（DSO） \\
CB-VIO \textsuperscript{[684]} & 基于模型 & 间接 & Individual events & 是 & 6 自由度 & 自然场景 & E+F+I & 特征跟踪器和 VIO 后端 \\
Wang \textsuperscript{[1154]} & 基于模型 & 间接 & Binary images & 是 & 6 自由度 & 自然场景 & 2E & 特征匹配 \\
El Moudni \textsuperscript{[303]} & 基于模型 & 直接 & Time Surfaces (TS) & 否 & 6 自由度 & 自然场景 & 2E & 使用 ESVO 跟踪器和 EMVS 建图器 \\
ESVIO \textsuperscript{[183]} & 基于模型 & 间接 & Time surfaces (TS) & 是 & 6 自由度 & 自然场景 & 2E+2F+I & 基于 TS 的特征跟踪 \\
ESVIO-direct \textsuperscript{[685]} & 基于模型 & 直接 & Time surfaces (TS) & 是 & 6 自由度 & 自然场景 & 2E+I & ESVO 的扩展 \\
PL-EVIO \textsuperscript{[410]} & 基于模型 & 间接 & Time surfaces (TS) & 是 & 6 自由度 & 自然场景 & E+F+I & 点与线特征、滑动窗口 BA \\
CMax-SLAM \textsuperscript{[416]} & 基于模型 & 直接 & Individual Events & 是 & 旋转 & 自然场景 & E & 对比度最大化细化运动 \\
EVI-SAM \textsuperscript{[409]} & 基于模型 & 直接、间接 & Individual Events & 是 & 6 自由度 & 自然场景 & E+F+I & 稠密建图 \\
Zuo \textsuperscript{[1317]} & 基于模型 & 直接 & Individual Event & 否 & 6 自由度 & 自然场景 & E+D & 使用深度传感器增强事件 \\
DEVO \textsuperscript{[585]} & 基于深度学习 & 间接 & Event voxel grids & 是 & 6 自由度 & 自然场景 & E & DPVO \textsuperscript{[1084]} 的事件版本 \\
EMBA \textsuperscript{[417]} & 基于模型 & 直接 & Individual Events & 是 & 旋转 & 自然场景 & E & 细化运动和梯度地图 \\
EPBA \textsuperscript{[418]} & 基于模型 & 直接 & Individual Events & 是 & 旋转 & 自然场景 & E & 细化运动和强度地图 \\
ES-PTAM \textsuperscript{[381]} & 基于模型 & 直接 & Events（也作为帧） & 否 & 6 自由度 & 自然场景 & 2E & 使用 EVO 跟踪器和 EMVS 建图器 \\
ESVO2 \textsuperscript{[812]} & 基于模型 & 直接 & Time surfaces (TS) & 是 & 6 自由度 & 自然场景 & 2E+I & ESVO 的扩展 \\
DEIO \textsuperscript{[408]} & 基于深度学习 & 间接 & Event voxel grids & 是 & 6 自由度 & 自然场景 & E+I & DEVO 和 DPVO 的扩展 \\
\bottomrule
\end{longtable}

表 10.1 基于事件的视觉 SLAM 方法汇总，按时间顺序排列。列分别表示方法类型（基于模型或基于深度学习、直接或间接）、是否具有全局细化模块（即后端/BA）、可处理的相机运动类型（旋转、平面、6 自由度）和场景类型（高对比度黑白场景等），以及输入数据类型（事件相机、帧式相机、深度传感器、IMU 和地图）；其中“2E”表示双目事件（两个传感器）。

\subsection{模拟器}

目前有多种公开工具可生成高质量的合成事件相机数据。ESIM \textsuperscript{[917]} 是 \textsuperscript{[778]} 的演进版本；\textsuperscript{[778]} 是最早模拟事件相机工作原理的模拟器之一。更早的工作（如 \textsuperscript{[535]}）只是对连续两帧之间的差异进行阈值化，以生成类似事件相机输出的边缘图像。ESIM 将渲染引擎与事件模拟器紧密耦合，使后者能够根据视觉信号的动态自适应地渲染帧。

CARLA \textsuperscript{[279]} 中的事件相机模拟器将 ESIM 扩展到更多样、丰富且复杂的自动驾驶场景。在基于事件学习单目深度 \textsuperscript{[460]} 的背景下，CARLA 开发的事件传感器从模拟器渲染图像中计算逐像素亮度变化，以与 ESIM 相同的方式模拟事件相机。图 10.5 展示了 CARLA 模拟器生成的 RGB 图像和事件。

\begin{figure}[htbp]
\centering
\includegraphics[width=\linewidth,height=0.78\textheight,keepaspectratio]{4e7ca98446094e7873ed9e09fb222edf4f6e41a66b35fde68883d0b051c279cd.jpg}
\label{fig:chapter-10-4e7ca98446094e7873ed9e09fb222edf4f6e41a66b35fde68883d0b051c279cd}
\end{figure}

(a) CARLA 模拟器

\begin{figure}[htbp]
\centering
\includegraphics[width=\linewidth,height=0.78\textheight,keepaspectratio]{53a5128d92d349e03120369137fef2a830d8b128d56580ea228bb9c929e80c95.jpg}
\label{fig:chapter-10-53a5128d92d349e03120369137fef2a830d8b128d56580ea228bb9c929e80c95}
\end{figure}

(b) V2E 模拟器
图 10.5 事件相机模拟器：(a) 在 CARLA 中使用 ESIM 模拟器生成的 RGB 图像和事件 \textsuperscript{[279]}；(b) V2E 工具的详细数据处理步骤 \textsuperscript{[481]}。

受基于学习的方法需要大量事件数据训练、而事件数据又因传感器新颖而难以获得这一事实的推动，研究者开发了将传统相机录制的任意视频转换为合成事件数据的工具 Video to Events（Vid2E）\textsuperscript{[371]}。Vid2E 旨在缩小传统视觉与事件视觉公开数据集之间的差距，使面向真实事件数据训练的网络可以使用几乎无限的现有视频数据集。Vid2E 将 ESIM 与自适应视频帧插值相结合，解决事件数据稀缺瓶颈。ESIM 能够以任意时间分辨率自适应渲染虚拟场景，但视频序列通常只提供固定且较低时间分辨率的强度测量。Super SloMo \textsuperscript{[521]} 可重建任意时间分辨率的中间帧，然后应用事件相机模拟器 ESIM。中间帧数量需要谨慎选择：数量过少会导致亮度信号混叠，数量过多则增加计算负担。

事件相机模拟器的一个重要方面是准确建模噪声，以缩小算法迁移时的仿真到真实差距。例如，ESIM 及其衍生模拟器根据经验观测 \textsuperscript{[661]} 实现了简单噪声模型：事件相机的对比度阈值（(10.1) 中的 $C$）并非固定，而是服从正态分布。为模拟这一点，ESIM 在每个模拟步骤从正态分布 $\mathcal { N } ( C , \sigma _ { C } ^ { 2 } )$ 采样，其中噪声水平 $\sigma _ { C }$ 可调。此外，ESIM 允许分别设置正、负对比度阈值（$C^{+}$ 和 $C^{-}$），以更准确地模拟真实事件相机。模拟器还应考虑其他噪声效应，例如电子噪声或事件像素带宽有限导致的对比度阈值空间和时间变化。

Vid2E \textsuperscript{[371]} 建模理想事件相机的照明。V2E \textsuperscript{[481]} 则基于 DVS 电路提出更真实的事件相机噪声模拟器。V2E 是首个包含时间噪声、泄漏事件和有限强度相关带宽的事件相机模拟器，并采用与 Vid2E 相同的高斯阈值分布。V2E 朝着更真实的模拟器迈进一步，可生成覆盖多种照明条件的合成数据集，这对事件应用十分重要。与 Vid2E 类似，V2E 使用 Super SloMo \textsuperscript{[521]} 提高输入视频的时间分辨率。图 10.5 详细展示了 V2E 的架构。

模拟事件相机噪声是实现真实合成事件生成的挑战。EventGAN \textsuperscript{[1306]} 提出一种端到端深度学习方案模拟事件相机数据：利用已有的图像标注数据，根据一对时间相邻的图像帧模拟事件，并使用 U-Net \textsuperscript{[934]} 编码器--解码器网络。该方法在图像和事件对上训练神经网络，不采用直接数值误差损失，而使用对抗判别器损失和一对循环一致性损失。EventGAN 为每种极性生成三维时空体素网格，而不是单个事件集合；这种表示常作为人工神经网络的输入。

VISTA 2.0 \textsuperscript{[30]} 是集成多种传感器类型（包括 RGB 相机、LiDAR 和事件相机）的模拟器，用于自动驾驶车辆的策略学习。它使用高保真的真实世界数据模拟多样场景，包括不同天气、照明和道路条件。事件相机模拟器与带自适应采样的 ESIM 类似，使用人工神经网络估计连续两帧之间的双向光流。VISTA 2.0 面向感知到控制策略的训练；与仅使用真实世界训练数据相比，它表现出更强的鲁棒性和仿真到真实迁移能力，从而改善安全关键场景中的车辆控制。

Video to Continuous Events（V2CE）\textsuperscript{[1282]} 旨在生成比既有模拟器更真实的事件时间戳。Vid2E 和 V2E 在离散时间戳生成事件，而真实事件以连续时间方式产生。对于对时间戳分布敏感的任务，这一差异不可忽略；它会造成显著域偏移，从而限制合成事件的使用。V2CE 分两阶段运行：第一阶段使用有监督的三维 U-Net 编码器--解码器 ANN，从视频预测两个体素网格（每种极性一个，与 EventGAN 类似）；第二阶段从体素网格恢复精确事件时间戳。该方法迭代推断每个体素中的事件数量及其相对位置。V2CE 还表明，它能够在饱和光照区域以及理想传感器事件生成模型不成立的边缘处准确生成事件。

\subsection{数据集与基准}

自 ECDS \textsuperscript{[778]} 发布以来，专用于视觉里程计和 SLAM 的事件数据集数量显著增加（见表 10.2）。ECDS 是首个同时提供同步事件、IMU 和 6 自由度相机真值位姿的数据集。更早的数据集 \textsuperscript{[955]} 包含合成和真实事件，场景简单、视觉对比度高，且仅有纯旋转运动（3 自由度）；真值由 IMU 获得。另一项工作 \textsuperscript{[55]} 支持事件相机与帧式相机运动的 5 自由度比较，其真值来自云台编码器和地面机器人的轮式里程计，因而容易漂移。ECDS 包含手持相机在多种场景中的 6 自由度运动（低速和高速），并由运动捕获系统提供精确的相机真值位姿。数据集包含 11 个真实事件场景和另外 2 个合成事件场景；合成数据由后来发展为 ESIM \textsuperscript{[917]} 的第一版模拟器生成。

RPG 双目数据集 \textsuperscript{[1301]} 包含在办公室环境中使用双目 DAVIS \textsuperscript{[111]} 录制的 8 段手持序列，以及由模拟器 \textsuperscript{[778]} 生成的一段合成序列（包含处于不同深度的三个正视平面）。该数据集不提供深度真值，但具有运动捕获系统提供的精确位姿真值，是事件双目 SLAM 方法原型开发和评估的良好起点。

Multi Vehicle Stereo Event Camera Dataset（MVSEC）\textsuperscript{[1304]} 是首个在多种平台上提供深度真值的数据集。它采集室内和室外场景，覆盖不同照明水平和运动速度。平台包括手持装置、六旋翼、汽车和摩托车，均配备经过标定的 3D LiDAR、IMU 和标准帧式相机。MVSEC 可用于位姿估计、建图、障碍物规避和三维重建等任务，并通过集成 LiDAR 提供准确的深度和位姿真值。数据集包含长序列，可全面评估基于事件的算法。

UZH-FPV（第一人称视角）数据集 \textsuperscript{[257]} 专为推动自主无人机竞速研究而设计。它采用定制四旋翼、Qualcomm Flight Board 和安装在 Lumenier QAV-R 碳纤维机架上的 mDAVIS346 \textsuperscript{[1074]} 事件相机。记录覆盖不同速度和轨迹的室内外场景，对导航和状态估计提出挑战。该数据集支持 VIO、事件数据处理和高速场景下的实时无人机应用研究，已成为开发高速相机运动算法的重要资源，并用于会议和研讨会竞赛。

Event Camera Motion Segmentation Dataset（EV-IMO）\textsuperscript{[769]} 是首个专门用于室内独立运动物体（IMO）分割的事件数据集。它包含 32 分钟记录，使用运动捕获系统跟踪最多三个快速运动的 IMO，并提供逐像素运动掩膜以及自运动和深度真值。该数据集适用于机器人学，尤其是需要精确运动检测的受限场景（如目标跟踪和自主导航）。EV-IMO2 \textsuperscript{[126]} 在其基础上增加了更多序列、三种更高质量的事件相机和更复杂的场景，既是当前算法的挑战性基准，也是开发新方法（包括单目和双目事件 SLAM）的丰富训练集。

表 10.2 事件 SLAM 数据集汇总，按时间顺序排列。传感器数据采用与表 10.1 相同的记号。双目和多传感器数据集的进一步介绍见综述 \textsuperscript{[380]}。

\small
\setlength{\tabcolsep}{2pt}
\begin{longtable}{p{0.237\linewidth}p{0.237\linewidth}p{0.237\linewidth}p{0.237\linewidth}}
\toprule
数据集 & 平台 & 像素分辨率 & 传感器 \\
\midrule
\endfirsthead
\toprule
数据集 & 平台 & 像素分辨率 & 传感器 \\
\midrule
\endhead
ECDS \textsuperscript{[778]} & 手持 & 240 × 180 & E、F、I \\
RPG-stereo \textsuperscript{[1301]} & 手持 & 240 × 180 & 2E \\
MVSEC \textsuperscript{[1304]} & 手持、无人机、汽车、摩托车 & 346 × 240 & 2E、2F、I、LiDAR、GPS \\
UZH-FPV \textsuperscript{[257]} & 无人机 & 346 × 260 & E、F、I \\
EV-IMO \textsuperscript{[769]} & 手持 & 346 × 260 & E、F、I、深度 \\
EV-IMO2 \textsuperscript{[126]} & 手持 & 640 × 480 & 3E、F、I、深度 \\
DSEC \textsuperscript{[372]} & 汽车 & 640 × 480 & 2E、2F、LiDAR、GPS \\
TUM-VIE \textsuperscript{[584]} & 手持 & 1280 × 720 & 2E、2F、I \\
EDS \textsuperscript{[461]} & 手持 & 640 × 480 & E、F（RGB）、I \\
VECtor \textsuperscript{[363]} & 手持 & 640 × 480 & 2E、2F、RGB-D、I、LiDAR \\
M2DGR \textsuperscript{[1250]} & 地面机器人 & 640 × 480 & E、F、I、LiDAR、GPS、热成像 \\
ViViD++ \textsuperscript{[635]} & 手持、汽车 & 240 × 180、640 × 480 & E、F、RGB-D、热成像、LiDAR、GPS \\
FusionPortable \textsuperscript{[523]} & 手持、四足机器人 & 346 × 240 & 2E、2F、I、LiDAR、GPS \\
Stereo HKU-VIO \textsuperscript{[183]} & 手持 & 346 × 260 & 2E、2F、I \\
M3ED \textsuperscript{[163]} & 无人机、汽车、四足机器人 & 1280 × 720 & 2E、2F、I、LiDAR、GPS \\
CoSEC \textsuperscript{[861]} & 汽车 & 1280 × 720 & 2E、2F、I、LiDAR、GPS \\
\bottomrule
\end{longtable}

\begin{figure}[htbp]
\centering
\includegraphics[width=\linewidth,height=0.78\textheight,keepaspectratio]{1f06c772e4c615375e0cfdc49d460583ea8c2cd026670d5fb0ce699dc59a408a.jpg}
\caption{图 10.6 部分事件 SLAM 数据集的细节：(a) DSEC 数据集 \textsuperscript{[372]} 中安装在汽车顶部的传感器套件（©2021 IEEE）；(b) EDS 数据集 \textsuperscript{[461]} 中使传感器共享空间对齐视场的分光器细节（©2022 IEEE）。}
\label{fig:chapter-10-1f06c772e4c615375e0cfdc49d460583ea8c2cd026670d5fb0ce699dc59a408a}
\end{figure}

Stereo Event Camera Dataset for Driving Scenarios（DSEC）\textsuperscript{[372]} 是面向自动驾驶的大规模数据集，尤其支持通过融合事件和帧开发能够应对不利照明条件的鲁棒感知系统。平台采用多相机配置，包括两台 VGA 分辨率 Prophesee Gen 3.1 事件相机（基线 60 cm）和两台 RGB 相机（FLIR Blackfly S，基线 51 cm，见图 10.6），同时配备 Velodyne VLP-16 LiDAR 和 RTK GPS 以实现精确定位。数据在瑞士不同城乡环境采集，涵盖白天、夜间和直射阳光等多种照明条件，并提供用于双目匹配的深度真值图。DSEC 还提供光流和视差，用于在挑战性驾驶条件下基准测试算法；基准指标包括 N 像素视差误差、平均绝对误差（MAE）和均方根误差（RMSE），以评估融合高分辨率事件相机数据与 RGB 帧的算法。

TUM Stereo Visual-Inertial Event Dataset（TUM-VIE）\textsuperscript{[584]} 使用双目 Prophesee Gen4 事件相机（1 百万像素）以及同步的 200 Hz IMU 数据和 20 Hz 双目灰度帧。它包含手持和头戴装置在多种室内外环境中的序列，涵盖体育活动、HDR 场景和低照度条件等。TUM-VIE 用于促进 VIO、SLAM、三维重建和传感器融合研究，尤其关注传统方法可能失效的挑战条件，推动高分辨率事件感知算法的发展。

Event-Aided Direct Sparse Odometry（EDS）数据集 \textsuperscript{[461]} 包含高质量事件、彩色帧和 IMU 数据，用于单目 VIO 研究。数据使用定制分光器装置采集（见图 10.6），使 RGB 帧和事件在同一光轴上精确对齐，这在以往数据集中并不常见。记录场景包括自然室内环境，提供高分辨率、良好标定的数据，可用于光流估计、深度估计，以及各种运动和照明条件下的鲁棒视觉里程计。

Versatile Event-Centric（VECtor）基准数据集 \textsuperscript{[363]} 同样用于评估事件 SLAM 算法。采集平台搭载多样传感器套件，包括双目相机（事件式和帧式）、RGB-D 传感器、128 通道 LiDAR 和九轴 IMU，均安装在灵活的 3D 打印支架上。数据集包含小规模室内环境（如运动捕获场）和具有不同复杂度及照明条件的大规模室内环境。据称，该数据集提供 VGA 分辨率、同步的多环境数据，支持在静态和动态、低照度及 HDR 场景中可靠评估 SLAM 算法，是推进多传感器 SLAM 研究的综合资源。

Vision for Visibility Dataset（ViViD++）\textsuperscript{[635]} 由包含热成像相机的多传感器平台采集，支持研究能够处理低能见度、运动干扰和外观变化的 SLAM 算法，利用不同传感器的互补优势。FusionPortable 数据集 \textsuperscript{[523]} 包含一台四足机器人，在走廊、食堂、道路和花园等不同场景中运动，覆盖不同照明条件。

最后，Multi-robot, Multi-Sensor, Multi-Environment Event Dataset（M3ED）\textsuperscript{[163]}（非正式名称为 MVSEC 2.0）聚焦机器人应用中的高速动态运动。它组合 1 百万像素双目事件相机、灰度和彩色相机、64 线 LiDAR 和高质量 IMU，全部与 RTK 定位同步。不同于早期数据集，M3ED 提供来自多种平台、结构化和非结构化环境的异构数据，并带有位姿、深度和语义标签真值，是开发面向传统驾驶或室内应用之外动态环境的鲁棒事件感知算法的重要资源。

\subsection{指标}

理想情况下，SLAM 系统应分别刻画定位模块和建图模块的质量。然而，(i) 两个模块以交织方式运行（深度误差会影响相机位姿误差，反之亦然），(ii) 定位真值信息（6 自由度）比精确深度更紧凑且更容易获取，因此深度估计误差通常被包含在相机轨迹误差的评估中。

由于经典 SLAM 和事件 SLAM 都输出相机轨迹，事件 SLAM 沿用经典 SLAM 的性能评估协议。常用的两个指标是绝对轨迹误差（ATE）和相对位姿误差（RPE）\textsuperscript{[1045]}。ATE 评估相机位姿相对于固定世界坐标系的精度，可全面反映 VO 系统的长期性能。RPE 评估连续（即相邻）时间步之间的相对位姿，关注 VO 系统的局部一致性。ATE 中的平移误差（也称位置误差）是估计相机位置与真值位置之间的欧氏距离；旋转误差（或方向误差）由 SO(3) 中的测地距离确定。类似地，RPE 的平移和旋转部分在时间间隔内的相机位姿对之间计算。有些研究还计算相对于场景平均深度或总行程的位置误差，以确保误差对场景或轨迹尺度不变。

还可使用其他误差指标评估估计的平移向量和旋转矩阵 \textsuperscript{[1305]}，包括平均 RPE（ARPE）、平均相对旋转误差（ARRE）和平均端点误差（AEE）。具体而言，ARPE 和 AEE 测量两个平移向量之间的位置和方向差异，ARRE 计算两个旋转矩阵之间的测地距离。

除上述指标外，平均线速度和角速度误差也可用于评估相机位姿估计，尤其是事件相机能够借助事件估计突发、快速相机运动的场景。相机位姿随时间变化，并由相应的线速度和角速度决定。已有多个工具箱 \textsuperscript{[405, 1283]} 公开发布，用于促进研究可复现性并降低 SLAM 轨迹评估的复杂性。

如果单独评估深度估计，通常使用直到固定深度值的多个截断处的平均深度误差，以比较不同三维地图尺度下的方法。然而，包含深度真值的数据集并不多（见第 10.7 节）。首选指标是估计三维点到真值地图最近表面之间欧氏距离的均方根误差（RMSE）。文献还使用相对误差（REL）和完整度（恢复点的数量）等指标 \textsuperscript{[302, 460, 381]}。

\section{延伸阅读与近期趋势}

尽管基于事件的 SLAM 研究已取得长足进展，但由于该技术仍较新，许多开放问题和挑战依然存在。这些问题涉及：对于给定任务（如 SLAM），如何采集和处理视觉信息（硬件和软件），使系统在鲁棒性、延迟、效率和精度等方面达到或超过生物系统的表现。

传感器是异步的，但由于神经形态计算的进入门槛，文献中的大多数系统仍设计在串行（即冯·诺依曼）处理器上。与预期的全神经形态系统 \textsuperscript{[845]} 相比，这在效率（功耗）、延迟等方面并非最优；全神经形态系统将事件相机与异步（受大脑启发、基于脉冲）的处理器、控制器和执行器等配对。构建能够模仿动物高效处理方式、在环境中自主建图和定位的机器人，一直是研究界的长期愿望（可能应用于 AR/VR 中“始终开启”的 inside-out 跟踪等）。实现这一愿望需要以神经形态工程方式重新思考并协同设计传感器、处理器和算法 \textsuperscript{[245]}，这极具挑战性，需要广博而深入的专业知识以及多个学科的协调。

在不久的将来，人们将开发受两条视觉通路启发的新型混合传感器 \textsuperscript{[1238]}，提供空间和时间对齐、低延迟、HDR 及精细细节（高像素数）的数据。研究者也在探索模仿生物学的中央凹传感器 \textsuperscript{[317]}，以降低带宽需求。事件相机及其演进（例如像素处理器阵列、Aeveon 传感器等近传感器处理器），以及与其他传感器（帧式相机、结构光、LiDAR、雷达等）的组合和数据融合，仍有广阔的研究空间，可进一步提升 SLAM 性能。

\section{致谢}

作者感谢 Giovanni Ciofi 在本章准备工作中提供的支持。
